\section{吞吐量革命：合并哲学与熵管理}

\subsection{当吞吐量改变合并哲学}

随着 Codex 产出速度的提高，许多传统的工程规范变得\textbf{适得其反}。

\begin{importantbox}{高吞吐量环境下的合并哲学}
仓库运行在最小化的阻塞合并门控（blocking merge gates）模式下：
\begin{itemize}[nosep]
  \item Pull Request 是短周期的
  \item 测试的 flake（不稳定失败）通常通过后续运行解决，而非无限期阻塞
  \item 在 Agent 吞吐量远超人类注意力的系统中，\textbf{修正是廉价的，等待是昂贵的}
\end{itemize}
\end{importantbox}

\begin{warningbox}{警告：这在低吞吐量环境中是不负责任的}
文章明确指出，这种``快速合并、后续修正''的策略\textbf{在低吞吐量环境中是不负责任的}。它只有在以下条件同时满足时才是合理的权衡：
\begin{itemize}[nosep]
  \item Agent 产出速度远超人类审查速度
  \item 修正的成本极低（Agent 可以快速修复）
  \item 架构约束提供了安全网（严格的 Linter 和测试）
  \item 每个变更是隔离的（独立的 worktree 和环境）
\end{itemize}
不满足这些前提时，照搬这种做法会导致灾难。
\end{warningbox}

\subsection{Agent 产出的全面性}

当文章说``代码库由 Codex Agent 生成''时，指的是\textbf{代码库中的所有东西}：

\begin{table}[H]
\centering
\begin{tabular}{ll}
\toprule
\textbf{产出类别} & \textbf{说明} \\
\midrule
产品代码和测试 & 应用逻辑与对应的自动化测试 \\
CI 配置和发布工具 & 持续集成管线和部署工具链 \\
内部开发者工具 & 辅助团队日常工作的脚本和工具 \\
文档和设计历史 & 架构决策记录与设计文档 \\
评估框架（Eval Harness） & 用于评估 Agent 产出质量的框架 \\
Review 评论和回复 & PR 上的代码审查交互 \\
仓库管理脚本 & 管理仓库本身的自动化脚本 \\
生产环境 Dashboard 定义 & 监控面板的配置文件 \\
\bottomrule
\end{tabular}
\caption{Codex Agent 的全部产出类别}
\end{table}

人类始终保持在循环中，但工作在不同的抽象层次：优先排序工作、将用户反馈转化为验收标准、以及验证结果。

\subsection{不断增长的自主性}

随着越来越多的开发循环被直接编码到系统中，仓库最近跨越了一个有意义的门槛：Codex 可以端到端地驱动一个新功能。给定单个 Prompt，Agent 现在可以：

\begin{enumerate}[nosep]
  \item 验证代码库的当前状态
  \item 复现一个被报告的 Bug
  \item 录制一段视频展示失败场景
  \item 实现修复
  \item 通过驱动应用验证修复
  \item 录制第二段视频展示修复后的结果
  \item 开启一个 Pull Request
  \item 响应 Agent 和人类的反馈
  \item 检测并修复构建失败
  \item 仅在需要人类判断时才升级
  \item 合并变更
\end{enumerate}

\begin{knowledgebox}{注意范围限定}
文章特别强调，这种高度自主的行为\textbf{严重依赖于该仓库的特定结构和工具链}，不应假设它能在没有类似投入的情况下泛化到其他项目——至少目前还不能。这是一个关于``什么是可能的''的演示，而非``所有人都应该这样做''的建议。
\end{knowledgebox}

\subsection{熵与``垃圾回收''}

完全 Agent 自主也引入了新问题。Codex 会\textbf{复制仓库中已存在的模式}——包括不均匀或次优的模式。随着时间推移，这不可避免地导致漂移。

最初，团队手动处理这个问题：\textbf{每周五花 20\% 的时间清理``AI 垃圾''（AI slop）}。不出所料，这种方式无法扩展。

\begin{importantbox}{Golden Principles + 自动化垃圾回收}
团队转而采用``黄金原则 + 自动清理''的策略：
\begin{enumerate}[nosep]
  \item 将``黄金原则''（golden principles）编码到仓库中——这些是保持代码库对未来 Agent 运行可读和一致的主观但机械化的规则
  \item 定期运行后台 Codex 任务，扫描偏差、更新质量评分、开启针对性的重构 PR
  \item 大多数这类 PR 可以在一分钟内审查并自动合并
\end{enumerate}
技术债务就像高利贷：\textbf{以小额持续还款几乎总比让它复利累积后痛苦地一次性偿还要好}。
\end{importantbox}

黄金原则的两个具体示例：
\begin{itemize}[nosep]
  \item \textbf{偏好共享工具包}：使用共享 utility package 而非手写 helper，以集中不变量
  \item \textbf{不做 YOLO 式数据探测}：在边界验证数据或依赖类型化 SDK，防止 Agent 基于猜测的数据形状构建逻辑
\end{itemize}

\subsection{本章小结}

高吞吐量环境下，``修正廉价、等待昂贵''成为合并哲学的核心。但这需要严格的架构约束作为安全网。Agent 会复制已有模式（包括坏模式），因此需要``垃圾回收''机制：将品味编码为黄金原则，通过自动化 Agent 持续扫描和修复偏差，实现技术债务的持续小额偿还。

\newpage

